\section{Weak Formulation}	\label{cha:WeakForm}
Consider again \cref{prb:Normalized_Problem}, i.e.,
\begin{subequations}
	\begin{equation}	\label{eqn:FEM_Problem}
		-\frac{d}{dx} \biggl[ \gls{q} \glsentryname{uPrime} \biggr] = \gls{f}, \quad x \in \left( 0, 1 \right),
	\end{equation}
	with boundary conditions
	\begin{align}
		\gls{u} \bigg \vert_{x = 0} &= \delta,  \label{eqn:FEM_Left_Boundary}  \\
		\biggl[ q \left( x \right) \frac{du}{dx} \biggr]_{x = 1} &= P. \label{eqn:FEM_Right_Boundary}
	\end{align}
\end{subequations}
Now introduce a \gls{r} by
\begin{equation}	\label{eqn:Residual_Function}
	r = \gls{r} = \frac{d}{dx} \left[ \gls{q} \glsentryname{uPrime} \right] + \gls{f} = 0,
\end{equation}
Next multiply both sides of \cref{eqn:Residual_Function} by arbitrary functions
\glsname{wi}\footnote{These are called \glspl{wi} and will be shortly specified.} and integrate the product over
$x \in \left[0, 1\right]$. Thus
\begin{subequations}
	\begin{align}
		\int_{x = 0}^{1} \gls{r} \gls{wi} \dd x &= 0,	\\
		\int_{x = 0}^{1} \frac{d}{dx} \left[ \gls{q} \frac{du}{dx} \right] \gls{wi} \dd x &+
		\int_{x = 0}^{1} \gls{f} \gls{wi} \dd x = 0. \label{eqn:FEM_Rewritten}
	\end{align}
\end{subequations}
%
Applying integration by parts to \cref{eqn:FEM_Rewritten} we get
\begin{equation}  \label{eqn:FEM2_Rewritten}
	\left[ \gls{q} \glsentryname{uPrime} \gls{wi} \right]_{x = 0}^{1}
	- \int_{x = 0}^{1} \gls{q} \frac{ du }{dx} \frac{d \glssymbol{wi}}{dx} \dd x = \int_{x = 0}^{1} \gls{f} \gls{wi} \dd x.
\end{equation}
Since the \gls{normalCondition} is only given on the \gls{rightBoundary}, i.e., at $x = 1$, we choose weight function such that
\begin{equation}
	\gls{wi} \Big \vert_{x = 0 } = 0.  \label{eqn:Weight_Condition}
\end{equation}
Now, insert the \gls{boundaryCondition} and the weight function property, i.e., \cref{eqn:FEM_Right_Boundary,eqn:Weight_Condition},
into \cref{eqn:FEM2_Rewritten} and rearrange the terms,
\begin{equation} \label{eqn:Weak_Formulation}
	\int_{x = 0}^{1} \gls{q} \glsentryname{uPrime} \frac{ d \glssymbol{wi} }{dx} \dd x = P  \gls{wi} \Big \vert_{x = 0} + \int_{x = 0}^{1} \gls{f} \gls{wi} \dd x.
\end{equation}
%
To summarize the \gls{weakForm} we need a couple of definitions. \\
\begin{definition} \label{def:Cont_of_Derivatives}
	Let $n \geq -1$ denote an integer. A general function of a single variable \gls{g} is said to support \gls{ordernCont},
	if and only if the function and its first $n$ derivatives exist and are continuous.
	Integer \glssymbol{orderOfCont} is then called \gls{orderOfCont}.
\end{definition}
Following terminology of \cref{def:Cont_of_Derivatives}, a single-variable function is known as:
\begin{itemize}
	\item discontinuous, if its \gls{orderOfCont} is $\glssymbol{orderOfCont} = -1$.
	\item infinitely differentiable or smooth, if its \gls{orderOfCont} is $\glssymbol{orderOfCont} = \infty$.
		Such a function supports \gls{orderInfCont}.
	\item continuous but not differentiable, if its \gls{orderOfCont} is $\glssymbol{orderOfCont} = 0$.
		Such a function supports \glsentryname{order0Cont}.
	\item up to one-time differentiable, if its \gls{orderOfCont} is $\glssymbol{orderOfCont} = 1$.
		Such a function supports \glsentryname{order1Cont}.
	\item up to two-times differentiable, if its \gls{orderOfCont} is $\glssymbol{orderOfCont} = 2$.
		Such a function supports \glsentryname{order2Cont}, and so on.
\end{itemize}
%
\begin{remark}
	Comparing left hand sides of \cref{eqn:FEM_Rewritten,eqn:Weak_Formulation} we observe that the second derivative involved in
	\cref{eqn:FEM_Rewritten} is replaced by the product of two first derivatives in \cref{eqn:Weak_Formulation}. The original
	requirement of \glsentryname{order1Cont} on \gls{u} is thus replaced by a weaker requirement on \gls{u} and \gls{wi},
	\glsentryname{order0Cont}. This is the reason why \cref{eqn:Weak_Formulation} is called a \gls{weakForm}.
\end{remark}
We conclude that \cref{prb:Normalized_Weak_Problem} below is a \gls{weakForm} of \cref{prb:Normalized_Problem}.
\begin{problem}	\label{prb:Normalized_Weak_Problem}
	Find a function $\gls{u}$ supporting \glsentrysymbol{order0Cont}-continuity such that
	\begin{subequations}
		\begin{align}	\label{eqn:Normalized_Weak_Diff_Eq}
			\int_{x = 0}^{1} \frac{d}{dx} \biggl[ \gls{q} \frac{du}{dx} \biggr] \gls{wi} \dd x
			+\int_{x = 0}^{1} \gls{f} \gls{wi} \dd x &= 0, & x \in \left( 0, 1 \right),
		\end{align}
		with boundary conditions
		\begin{align}
			\gls{u} \bigg \vert_{x = 0} &= \delta,  \label{eqn:Normalized_Weak_Left_Boundary} \\
			\Biggl[ \gls{q} \glsentryname{uPrime} \Biggr]_{x = 1} &= P. \label{eqn:Normalized_Weak_Right_Boundary}
		\end{align}
	\end{subequations}	
\end{problem}
